\chapter{Overview}
\label{ch:overview}

\haira{} is a statically-typed, compiled programming language that brings AI-assisted code generation into the compilation pipeline. It compiles to native binaries via Cranelift, with no runtime dependencies.

\section{Hello, World}

\begin{lstlisting}
import "io"

fn main() {
    io.println("Hello, World!")
}
\end{lstlisting}

To compile and run:

\begin{lstlisting}[language=bash,numbers=none]
$ haira build hello.haira
$ ./hello
Hello, World!
\end{lstlisting}

\section{Language Characteristics}

\subsection{Explicit Imports}

All dependencies are explicitly declared at the top of each file:

\begin{lstlisting}
import "io"
import "json"
import db from "haira/postgres"
\end{lstlisting}

\subsection{Go-Style Error Handling}

Functions that can fail return a tuple of result and error:

\begin{lstlisting}
import "io"
import "fs"

fn main() {
    content, err = fs.read_file("config.json")
    if err != nil {
        io.println("Error: " + err.message)
        return
    }
    io.println(content)
}
\end{lstlisting}

\subsection{Explicit AI Blocks}

AI-generated code is explicit, using the \keyword{ai} keyword:

\begin{lstlisting}
import "io"

struct User {
    name: string
    email: string
    posts: []Post
}

ai summarize_user(user: User) -> string {
    Generate a brief, human-readable summary of this user
    including their name and post count.
}

fn main() {
    user = User{
        name: "Alice",
        email: "alice@example.com",
        posts: []
    }
    io.println(summarize_user(user))
}
\end{lstlisting}

\begin{note}
AI blocks are processed at \textbf{compile time}. The generated code is cached and included in the binary. There are no AI dependencies at runtime.
\end{note}

\subsection{Local-First AI}

AI code generation uses local models by default:

\begin{enumerate}
    \item \textbf{llama.cpp} -- Primary backend, runs local GGUF models
    \item \textbf{Ollama} -- Secondary, connects to local Ollama server
    \item \textbf{Cloud APIs} -- Optional fallback (Anthropic, OpenAI)
\end{enumerate}

Configuration in \code{haira.toml}:

\begin{lstlisting}[language={}]
[ai]
backend = "llama.cpp"
model_path = "~/.haira/models/codellama-7b.gguf"
\end{lstlisting}

\section{Type System}

\haira{} uses structural typing with type inference for local variables:

\begin{lstlisting}
// Type inferred as int
count = 42

// Explicit type annotation
name: string = "Haira"

// Struct definition
struct Point {
    x: float
    y: float
}

// Option type (nullable)
maybe_value: int? = nil

// Arrays and maps
numbers: []int = [1, 2, 3]
scores: [string:int] = ["alice": 100, "bob": 95]
\end{lstlisting}

\section{Control Flow}

\begin{lstlisting}
import "io"

fn main() {
    // If-else
    x = 10
    if x > 5 {
        io.println("Large")
    } else {
        io.println("Small")
    }

    // For loop with range
    for i in 0..5 {
        io.println(i)
    }

    // While loop
    n = 0
    while n < 3 {
        io.println(n)
        n = n + 1
    }

    // Match expression
    status = 200
    match status {
        200 => io.println("OK")
        404 => io.println("Not Found")
        500 => io.println("Server Error")
        _ => io.println("Unknown")
    }
}
\end{lstlisting}

\section{Functions}

Functions are declared with \keyword{fn}:

\begin{lstlisting}
fn add(a: int, b: int) -> int {
    return a + b
}

// Multiple return values
fn divide(a: int, b: int) -> (int, Error?) {
    if b == 0 {
        return 0, Error{message: "division by zero"}
    }
    return a / b, nil
}

// Closures
fn make_adder(n: int) -> fn(int) -> int {
    return fn(x: int) -> int {
        return x + n
    }
}
\end{lstlisting}

\section{Concurrency}

\haira{} provides Go-style concurrency with \keyword{spawn} and channels:

\begin{lstlisting}
import "io"

fn main() {
    ch = chan<int>(10)  // Buffered channel

    spawn {
        for i in 0..5 {
            ch <- i
        }
        close(ch)
    }

    for value in ch {
        io.println(value)
    }
}
\end{lstlisting}

\section{Pipe Operator}

The pipe operator \code{|} enables functional composition:

\begin{lstlisting}
import "io"
import "string"
import "array"

fn main() {
    result = "  hello, world  "
        | string.trim
        | string.to_upper
        | string.split(", ")
        | array.first

    io.println(result)  // "HELLO"
}
\end{lstlisting}

\section{A Complete Example}

\begin{lstlisting}
import "io"
import "http"
import "json"

struct User {
    id: int
    name: string
    email: string
}

ai validate_email(email: string) -> (bool, string?) {
    Validate the email format and return whether it's valid.
    If invalid, return an error message explaining why.
}

fn create_user(name: string, email: string) -> (User, Error?) {
    valid, reason = validate_email(email)
    if not valid {
        return User{}, Error{message: reason or "Invalid email"}
    }

    user = User{
        id: generate_id(),
        name: name,
        email: email
    }

    return user, nil
}

fn main() {
    user, err = create_user("Alice", "alice@example.com")
    if err != nil {
        io.println("Failed: " + err.message)
        return
    }

    io.println("Created user: " + user.name)
}

fn generate_id() -> int {
    // Simple ID generation
    return 1
}
\end{lstlisting}

\section{Implementation Status}

This specification documents the complete \haira{} language. Implementation proceeds incrementally:

\begin{table}[h]
\centering
\begin{tabular}{llc}
\toprule
\textbf{Feature} & \textbf{Chapter} & \textbf{Target Release} \\
\midrule
Lexer \& Parser & 2--3 & v0.1 \\
Expressions \& Statements & 4--5 & v0.2 \\
Functions & 6 & v0.3 \\
Error Handling & 7 & v0.3 \\
Concurrency & 8 & v0.4 \\
Modules \& Stdlib & 9--10 & v0.5 \\
Advanced Features & 11 & v0.6 \\
AI Integration & 12--13 & v1.0 \\
\bottomrule
\end{tabular}
\caption{Implementation roadmap by chapter}
\end{table}

\begin{implementation}
Features not yet implemented will produce clear error messages indicating the target release version.
\end{implementation}
